\section{Algebraic boundary and novelty discipline}
\label{sec:boundary}

Four mature theories meet in SD-C17.  Keeping their data types separate is
more important than expanding the bibliography.

\subsection{Weighted symbolic determinants and necklaces}

For a one-vertex shift with finitely or absolutely summably many weighted
edges, the vertex adjacency is the sum of the edge weights.  Its determinant
is therefore $1-F$, and
\[
  -\log(1-F)=\sum_{n\ge1}\frac{F^n}{n}
  \tag{2.1}
\]
regroups into primitive cyclic words and their repetitions.  This is the
weighted one-vertex specialization of the symbolic determinant framework of
\citet{bowenlanford1970}.  Necklaces, primitive roots, and their Witt-type
algebras are classical \citep{metropolisrota1983}.  We use only the elementary
primitive-root decomposition and prove the required finite counts directly.

Equation (2.1) already warns against confusing three objects:

\begin{itemize}
\item a signed edge in the alphabet;
\item a primitive cyclic word made of those edges;
\item the $r$th temporal repetition of that primitive word.
\end{itemize}

The determinant aggregates all three.  A primitive-orbit claim must recover
them before aggregation.

\subsection{Bar and Koszul reductions}

For an augmented algebra, the bar resolution is canonical but large.
Koszul algebras admit smaller resolutions with exterior-shaped generators
\citep{priddy1970}; modern Koszul duality places this pattern in a broader
operadic setting \citep{ginzburgkapranov1994}.  Algebraic discrete Morse
theory constructs chain-homotopy reductions of bar and related complexes
\citep{skoldberg2006,jollenbeckwelker2009}.  If a group acts, the matching
must satisfy additional equivariance conditions; arbitrary order-based
matchings need not descend equivariantly \citep{freij2009}.

These results authorize a genuine chain reduction only when chain groups,
differentials, matchings, and induced symmetries are specified.  They do not
authorize replacing a scalar coefficient $(-1)^d$ by a supertrace parity
after a determinant has been computed.  SD-C17 is intentionally the scalar
decategorification, and \Cref{sec:homological} tests whether that scalar object
can be lifted back without changing its temporal ledger.

\subsection{Hochschild, Harrison, and indecomposables}

For a smooth polynomial algebra in characteristic zero, the
Hochschild--Kostant--Rosenberg theorem identifies Hochschild homology with
differential forms \citep{hkr1962}.  Mixed forms survive in exterior degrees
at least two.  Harrison theory isolates the commutative indecomposable
component \citep{barr1968}; André--Quillen theory expresses the same smooth
direction through the cotangent complex \citep{quillen1970}.  Hodge-type
decompositions make explicit that the Harrison component is one summand or
quotient of the fuller Hochschild information, not its entirety
\citep{gerstenhaberschack1987,loday1998}.

This distinction blocks a tempting shortcut.  An atom-only Harrison or
cotangent sector may be mathematically natural, but projecting to it deletes
mixed cyclic recurrence.  It cannot be presented as a trace-preserving
equivalence of the original scalar subset shift.

\subsection{Exact novelty boundary}

The following are not new here: inclusion--exclusion, the Euler product in
its convergence half-plane, the Koszul resolution, the exterior computation
of Tor, HKR, Harrison/AQ decomposition, equivariant Morse theory, and
necklace enumeration.  We claim only the following bounded contribution:

\begin{quote}
For the source-locked tensor-atom subset shift, the $p^2q^2$ power ledger,
the $pqr$ $S_3$ character, and the scalar/supertrace firewall jointly rule
out a primitive-local, power-compatible, atom-permutation-natural Koszul lift.
\end{quote}

No priority claim is made for the individual combinatorial identities.  The
value is the obstruction package and the decision it forces for SD-C17.
